\begin{theorem}[全整数预算成功曲线的完全充分性]\label{thm:conclusion-budget-curve-exact-tail-difference-reconstruction}
固定分辨率 \(m\)。定义全整数预算成功曲线
\[
\mathrm{Succ}_m(s):=
2^{-m}\sum_{x\in X_m}\min\bigl(d_m(x),s\bigr),
\qquad
s=0,1,\dots,2^m,
\]
并记
\[
C_m(s):=2^m\mathrm{Succ}_m(s),
\qquad
\Delta_m(s):=C_m(s)-C_m(s-1)
\qquad (1\le s\le 2^m),
\]
以及纤维多重度直方图
\[
M_m(k):=\#\{x\in X_m:\ d_m(x)=k\}.
\]
则有精确公式
\[
\Delta_m(s)=\#\{x\in X_m:\ d_m(x)\ge s\},
\]
\[
M_m(k)=\Delta_m(k)-\Delta_m(k+1).
\]
因此整条成功曲线 \(s\mapsto \mathrm{Succ}_m(s)\) 唯一决定纤维多重度多重集 \(\{d_m(x):x\in X_m\}\)。
\end{theorem}

\begin{proof}
对任意固定整数 \(d\ge 0\) 与 \(s\ge 1\)，有恒等式
\[
\min(d,s)-\min(d,s-1)=\mathbf 1_{\{d\ge s\}}.
\]
将其应用到每条纤维多重度 \(d_m(x)\) 并对 \(x\in X_m\) 求和，得到
\[
\Delta_m(s)
=
\sum_{x\in X_m}
\Bigl(\min(d_m(x),s)-\min(d_m(x),s-1)\Bigr)
=
\#\{x\in X_m:\ d_m(x)\ge s\}.
\]
再用尾计数差分恒等式
\[
\#\{d_m\ge k\}-\#\{d_m\ge k+1\}
=
\#\{d_m=k\}
\]
即得
\[
M_m(k)=\Delta_m(k)-\Delta_m(k+1).
\]
于是 \(\mathrm{Succ}_m\) 唯一决定全部 \(M_m(k)\)，进而唯一决定纤维多重度多重集。证毕。
\end{proof}

\begin{corollary}[预算层即 Mellin--Rényi 因子化层]\label{cor:conclusion-budget-curve-mellin-renyi-factorization-layer}
设 \((\mathrm{Succ}_m)_{m\ge 1}\) 为全部分辨率上的整数预算成功曲线塔。则对任何仅依赖纤维多重度多重集的观测量，它都是完全充分统计量。特别地，下列对象都经由预算曲线塔唯一因子化：
\[
S_q(m)=\sum_{k\ge 1}k^q M_m(k)
\qquad (q>0),
\]
\[
C_m^{\mathrm{cont}}(T)=\sum_{k\ge 1}\min(k,T)M_m(k),
\]
\[
N_m(t)=\sum_{k\ge \lceil t\rceil}M_m(k).
\]
因而全部 Mellin 矩、连续容量曲线、尾谱、以及由这些对象诱导的压力谱 \(P_q\)、冻结阈值、基态简并指数 \(g_\ast\) 与容量指数包络，都经由预算曲线塔唯一决定。
\end{corollary}

\begin{proof}
由定理 \ref{thm:conclusion-budget-curve-exact-tail-difference-reconstruction}，预算曲线塔唯一决定全部直方图 \(M_m(k)\)。而题中各式都是对 \(M_m(k)\) 的显式函数，因此均由预算曲线塔唯一决定。项目文稿中压力、冻结阈值、基态简并指数与容量指数包络的定义只依赖于这些由直方图导出的量，故结论成立。证毕。
\end{proof}

\begin{proposition}[可逆容量指数的普适折角]\label{prop:conclusion-capacity-exponent-universal-fold-angle}
记
\[
\Gamma(\beta):=
\limsup_{m\to\infty}\frac1m\log C_m(\lfloor \beta m\rfloor),
\qquad
\beta\ge 0,
\]
并定义
\[
M_m^{\max}:=\max_{x\in X_m}d_m(x),
\qquad
\alpha_\ast:=\lim_{m\to\infty}\frac1m\log M_m^{\max},
\qquad
g_\ast:=\lim_{m\to\infty}\frac1m\log K_m.
\]
则项目文稿给出的普适下支满足
\[
\Gamma(\beta)\ge g_\ast+\min\{\alpha_\ast,\beta\log 2\}.
\]
若记
\[
\beta_c:=\frac{\alpha_\ast}{\log 2},
\]
则该下支
\[
\underline\Gamma(\beta):=g_\ast+\min\{\alpha_\ast,\beta\log 2\}
\]
在 \(\beta_c\) 处具有不可消除的折角：
\[
\underline\Gamma(\beta)=
\begin{cases}
g_\ast+\beta\log 2,& \beta\le \beta_c,\\[1mm]
g_\ast+\alpha_\ast,& \beta\ge \beta_c.
\end{cases}
\]
其左斜率为 \(\log 2\)，右斜率为 \(0\)。
\end{proposition}

\begin{proof}
这是项目文稿中容量指数双侧夹逼之下支的直接改写。由
\[
\beta\log 2=\alpha_\ast
\]
的唯一交点恰为
\[
\beta_c=\frac{\alpha_\ast}{\log 2},
\]
故 \(\underline\Gamma\) 的显式分段式立即成立，左右导数分别为 \(\log 2\) 与 \(0\)。证毕。
\end{proof}

\begin{theorem}[有界实现切片上的 bulk--boundary 有限完备审计]\label{thm:conclusion-bulk-boundary-joint-finite-complete-audit}
固定分辨率 \(m\)，并假设对象属于定理 \ref{thm:conclusion-toeplitz-endpoint-atom-dyadic-joint-verification} 的有界实现--端点谱隙切片。定义审计对象
\[
\mathcal A_{m,D,\varepsilon}
:=
\Bigl(
(\mathrm{Succ}_m(s))_{s=0}^{2^m},
\ (T_{2^k})_{k=0}^{K(D,\varepsilon)},
\ a_{2^{K(D,\varepsilon)}}
\Bigr).
\]
则 \(\mathcal A_{m,D,\varepsilon}\) 同时完成下列三件互异任务：
\begin{enumerate}
  \item \textbf{bulk 重构：}唯一恢复纤维多重度直方图 \(M_m(k)\)；
  \item \textbf{global feasibility：}判定全部 Toeplitz--PSD 约束；
  \item \textbf{boundary extraction：}以误差不超过 \(\varepsilon\) 的精度恢复端点原子质量 \(\mu(\{-1\})\)。
\end{enumerate}
\end{theorem}

\begin{proof}
由定理 \ref{thm:conclusion-budget-curve-exact-tail-difference-reconstruction}，第一分量 \((\mathrm{Succ}_m(s))_{s=0}^{2^m}\) 唯一恢复全部直方图 \(M_m(k)\)，从而完成 bulk 重构。由定理 \ref{thm:conclusion-toeplitz-endpoint-atom-dyadic-joint-verification}，第二分量 \((T_{2^k})_{k=0}^{K(D,\varepsilon)}\) 判定全层级 Toeplitz--PSD，可完成 global feasibility；第三分量 \(a_{2^{K(D,\varepsilon)}}\) 以误差不超过 \(\varepsilon\) 恢复 \(\mu(\{-1\})\)，完成 boundary extraction。三者合并即得结论。证毕。
\end{proof}
